% ============================================================
% CAPITOLO 19 - LE NOTE ESTRANEE ALL'ARMONIA
% ============================================================
\chapter{Le note estranee all'armonia}
\label{chap:noteestranee}

Le note che appartengono al tessuto melodico ma non all'accordo sottostante si dicono \textbf{estranee} (o non armoniche). Arricchiscono la trama e rendono le melodie più fluide e cantabili.

\begin{definizione}[Nota estranea all'armonia]
    \label{def:nota_estranea}
    Una nota è \textbf{estranea} (o non armonica) quando non fa parte dell'accordo che suona in quel momento. Crea una dissonanza temporanea che risolve su una nota reale dell'accordo.
\end{definizione}

\section{Nota di passaggio}
\label{sec:nota_passaggio}

\begin{definizione}[Nota di passaggio]
    È una nota estranea che occupa per grado congiunto l'intervallo tra due note reali dell'accordo, consecutive o appartenenti ad accordi diversi.
\end{definizione}

\begin{regola}[Uso della nota di passaggio]
    La nota di passaggio cade di norma sul tempo debole (o sulla parte debole del tempo). Se cade sul tempo forte assume un carattere più teso, simile all'appoggiatura. Può essere \textbf{diatonica} (nella scala) o \textbf{cromatica} (alterata).
\end{regola}

\begin{esempio}[Nota di passaggio in Do maggiore]
    Melodia Do-Re-Mi sopra un accordo di Do maggiore: il Re è la nota di passaggio.
    \musica{es_cap11_01}
\end{esempio}

\begin{esempio}[Nota di passaggio a quattro parti]
    \musica{es_cap11_03}
\end{esempio}

\section{Nota di volta}
\label{sec:nota_volta}

\begin{definizione}[Nota di volta]
    È una nota che si allontana dalla nota reale per grado congiunto (verso l'alto o verso il basso) e vi ritorna subito.
\end{definizione}

\begin{regola}
    La nota di volta cade di solito sul tempo debole. La nota di ritorno è identica alla nota reale di partenza.
\end{regola}

\section{Anticipazione}
\label{sec:anticipazione}

\begin{definizione}[Anticipazione]
    È una nota che appartiene all'accordo \textit{successivo} ma compare prima del cambio armonico, sul tempo debole.
\end{definizione}

\begin{regola}
    A differenza del ritardo, l'anticipazione non è preparata e non risolve: \textit{è} già la risoluzione anticipata.
\end{regola}

\begin{esempio}[Anticipazione nella cadenza V-I]
    La tonica del soprano compare sul tempo debole, mentre il basso tiene ancora la dominante, prima del cambio d'accordo.
    \musica{es_cap11_02}
\end{esempio}

\section{Appoggiatura}
\label{sec:appoggiatura}

\begin{definizione}[Appoggiatura]
    È una nota estranea posta sul tempo \textbf{forte}, raggiunta senza preparazione (di solito per salto, talvolta per grado) e risolta per grado congiunto discendente (più raramente ascendente).
\end{definizione}

\begin{regola}[Appoggiatura e ritardo]
    A differenza del ritardo (Capitolo \ref{chap:ritardi}), l'appoggiatura non è preparata da una nota tenuta. Storicamente era spesso scritta come abbellimento, ma ha il peso metrico di una nota reale.
\end{regola}

\section{Nota di sfuggita}
\label{sec:sfuggita}

\begin{definizione}[Nota di sfuggita (échappée)]
    È una nota estranea raggiunta per grado congiunto e abbandonata per salto in direzione opposta. È comune nel soprano nei movimenti melodici veloci.
\end{definizione}

\section{Bordone e nota pedale}
\label{sec:pedale}

\begin{definizione}[Nota pedale]
    È una nota tenuta, di solito al basso o al soprano, mentre l'armonia cambia sopra o sotto di essa.
\end{definizione}

\begin{regola}
    La nota pedale è consonante con l'accordo all'inizio e alla fine del passaggio. Nel mezzo sono ammesse dissonanze, che accumulano tensione.
\end{regola}
